% ===== 原文 pp.142--149（PDF p.152--159）：附录 C 相伴纤维丛与等变上同调 =====
% ===== 本文件 appC-02.tex：印刷 pp.146--149（PDF p.156--159） =====

% [原文 p.146]
\begin{definition}
由底层主丛 $P$ 中的联络 $\omega$ 诱导的截面 $\bm{\Phi}$ 的协变导数，是元素 $\dd_\omega\bm{\Phi}\in\Omega^1(M,\bm{\Phi}^*VE)$，它由前推与垂直投影的复合给出：
\[ \dd_\omega\bm{\Phi} = \omega^V_E\circ\dd\bm{\Phi}: TM\to TE\to VE . \]
\end{definition}

% [Tutor Note: p.146 原文把 principal bundle 误拼为 “principle bundle”（本页与 p.147 各出现一次）——应作 principal bundle（主丛），译文按正确含义处理。另核对 $\dd_\omega\bm{\Phi}$ 的取值空间：$\dd\bm{\Phi}$ 映到 $TE$，$\omega^V_E$ 再映到 $VE$，而 $\bm{\Phi}^*VE$ 是 $M$ 上的丛，故 $\dd_\omega\bm{\Phi}\in\Omega^1(M,\bm{\Phi}^*VE)$ 是 1-形式取值于该丛的截面。]

注意，我们用了 $P$ 上一个具体的联络选取来定义上面的复结构 $J_E$。因此，例如“给定截面 $\bm{\Phi}: M\to E$ 是复流形之间的全纯映射”这一条件将依赖于 $\omega$ 的选取。

% [Tutor Note: p.146 为何全纯性依赖 $\omega$：按 §C.1 的构造，$J_E$ 在 $VE$ 上等于 $J^V$（由 $J_X$ 拉回，与 $\omega$ 无关），而它在 $HE$ 上等于 $J^H$，且水平分解 $TE=HE\oplus VE$ 本身就由 $\omega$ 决定。核对：换一个联络 $\omega'$ 只改变子丛 $HE\subset TE$，而全纯性条件 $\bar\partial\bm{\Phi}=0$ 的 $(0,1)$ 部分含水平分量，故该条件随 $\omega$ 改变。]

\section{拉回 \texorpdfstring{$X$}{X} 上的等变上同调类}

关于相伴纤维丛的另一种互补观点，是把截面及其协变导数看作 $P$（而非 $M$）上的整体定义量：

\begin{theorem}
给定截面 $\bm{\Phi}: M\to E$，存在唯一的映射 $\varphi: P\to X$ 使得
\[ \bm{\Phi}(\pi(p)) = [p,\varphi(p)] \]
对一切 $p\in P$ 成立。该映射在如下意义下是 $G$-等变的：
\[ \varphi(p\cdot g) = g^{-1}\cdot\varphi(p) \]
对一切 $p\in P$ 与 $g\in G$ 成立。反之，一个 $G$-等变映射 $\varphi$ 定义了一个映射 $\bm{\Phi}$，因此我们可以等价地把相伴纤维丛的截面看成从 $P$ 到 $X$ 的 $G$-等变映射。
\end{theorem}

% [Tutor Note: p.146 等变条件的幂次可用一致性核对：$\bm{\Phi}(\pi(p\cdot g)) = \bm{\Phi}(\pi(p)) = [p,\varphi(p)]$，而同一个点在 $P$ 中以 $p\cdot g$ 表示时应写作 $[p\cdot g,\ \varphi(p\cdot g)] = [p,\ g\cdot\varphi(p\cdot g)]$。两式相等给出 $\varphi(p\cdot g)=g^{-1}\cdot\varphi(p)$；若误写成 $g$ 则左、右作用不匹配（$P$ 是右作用、$X$ 是左作用）。]

\begin{definition}
一个 $G$-等变映射 $\varphi: P\to X$ 关于 $P$ 上联络 $\omega$ 的协变导数，是元素 $\dd_\omega\varphi\in\Omega^1(P,\varphi^*TX)$，它由
\[ \dd_\omega\varphi \equiv \dd\varphi\circ\omega^H: TP\to TX \]
给出。它与上面定义的协变导数 $\dd_\omega\bm{\Phi}$ 由
\[ \dd_\omega\bm{\Phi}(\dd\pi(u_p)) = [u_p,\ \dd_\omega\varphi] \]
相联系，对一切 $u_p\in TP$ 成立。
\end{definition}

\begin{theorem}
协变导数 $\dd_\omega\varphi$ 由如下公式给出：
\[
    \dd_\omega\varphi|_p = \dd\varphi|_p + \omega^a|_p\, V_a|_{\varphi(p)}: T_pP\to T_{\varphi(p)}X
\]
\end{theorem}

% [原文 p.147]
其中 $V_a\in TX$ 是对应于李代数基元 $t_a$（$a=1,\dots,\dim G$）的 $X$ 上的向量场，而 $\omega\in\Omega^1(P,\fg)$ 是 $P$ 上的联络 1-形式。

\begin{proof}
我们写
\[
    \dd_\omega\varphi|_p = \dd\varphi\circ\omega^H|_p = (\dd\varphi-\dd\varphi\circ\omega^V)|_p = (\dd\varphi-\dd\varphi\circ\xi\circ\omega)|_p,
\]
其中 $\xi|_p: \fg\to T_pP$ 是对应于无穷小右 $G$-作用的映射。然而，由于 $\varphi$ 的 $G$-等变性，我们有
\[ (\dd\varphi\circ\xi)|_p: t_a \longmapsto -V_a|_{\varphi(p)}, \]
其中负号是因为 $G$ 从右边作用在 $P$ 上，而从左边作用在 $X$ 上。代入上式即证明该公式。
\end{proof}

% [Tutor Note: p.147 证明中 $(\dd\varphi\circ\xi)|_p$ 的负号可独立核对：把 $G$-等变性 $\varphi(p\cdot g)=g^{-1}\cdot\varphi(p)$ 在 $g=e$ 处对 $g$ 求导（沿 $t_a$ 方向），左边给出 $\dd\varphi(\xi(t_a))$，右边给出 $-V_a|_{\varphi(p)}$，即得该式；这也正是“$G$ 右作用于 $P$、左作用于 $X$”导致符号相反的地方。]

规范 $\sigma$-模型作用量中的拓扑项对应于靶流形 $X$ 的等变上同调类，下面我们就来说明这一点。我们选择在 $X$ 的等变上同调的 Cartan 模型中工作。$X$ 上的等变 $k$-形式是复形中一个 $k$ 次元素 $\eta$：
\[
    \eta\in\Omega^k_G(X) = S(\fg^*)\otimes\left[\Omega^\bullet(X)\right]^G\Big|_{\text{degree }k},
\]
其中 $S(\fg^*)$ 是 $\fg^*$ 上的对称代数，而 $\left[\Omega^\bullet(X)\right]^G$ 是被 $N$ 个李导数
\[ \cL_a \equiv L_{V_a} = \dd_X\circ\iota_{V_a} + \iota_{V_a}\circ\dd_X \]
零化的所有 $X$ 上形式构成的子空间。这是一个 $\mathbb{Z}$-分次向量空间，其中 $S(\fg^*)$ 的每个生成元 $\zeta^a$ 被形式地赋予次数二，$\Omega^j(X)$ 的每个元素被赋予次数 $j$。$\eta$ 的最一般形式是
\[
    \eta = \sum_{i_1=0}^\infty\sum_{i_2=0}^\infty\cdots\sum_{i_N=0}^\infty
    \eta_{(i_1,i_2,\dots,i_N)}(\zeta^1)^{i_1}(\zeta^2)^{i_2}\cdots(\zeta^N)^{i_N}
    \equiv \sum_I \eta_I\zeta^I,
\]
其中系数 $\eta_I=\eta_{(i_1,\dots,i_N)}$ 由 $N$ 元非负整数数组标号，并且
\[ \eta_I\in\left[\Omega^{k-2|I|}(X)\right]^G, \]
其中 $|I| = i_1+\cdots+i_N$。给定 $X$ 上这样一个等变 $k$-形式 $\eta$，可以用一个截面 $\varphi: P\to X$ 与 $P$ 上的一个联络 $\omega$ 把 $\eta$ “拉回”成 $P$ 上的普通 $k$-形式：

% [Tutor Note: p.147 次数核对：生成元 $\zeta^a$ 被赋予次数 2，故含 $|I|$ 个生成元的项里 $\eta_I$ 必须是次数 $k-2|I|$ 的形式，形式总次数才能为 $k$；这解释了为何 $\eta_I$ 只可能对 $|I|\leqslant k/2$ 非零（原文把求和上限写成 $\infty$，实际上有限项即止）。]

\begin{definition}
令 $\eta\in\Omega^k_G(X)$ 为如上等变 $k$-形式，$\varphi: P\to X$ 为某个截面所确定的等变映射，$\omega\in\Omega^1(P,\fg)$ 为 $P$ 上的联络，其曲率为 $F_\omega = \dd\omega+\half[\omega,\omega]\in\Omega^2(P,\fg)$。
% [原文 p.148]
三元组 $(\eta,\varphi,\omega)$ 如下定义一个普通 $k$-形式 $\eta(\varphi,\omega)\in\Omega^k(P)$：
\[
    \eta(\varphi,\omega) = \sum_I \left[(\dd_\omega\varphi)^*\eta_I\right]\wedge F^I_\omega .
\]
\end{definition}

也就是说，我们用协变导数 $\dd_\omega\varphi$ 把 $\eta$ 的系数拉回到 $P$ 上，并把 $N$ 个生成元 $\zeta^1,\dots,\zeta^N$ 换成 $N$ 个 2-形式 $F^1_\omega,\dots,F^N_\omega$。$k$-形式 $\eta(\varphi,\omega)$ 具有以下性质：

% [Tutor Note: p.148 此处 $(\dd_\omega\varphi)^*\eta_I$ 是 $X$ 上的形式沿映射 $\dd_\omega\varphi: TP\to TX$ 的拉回（不是沿 $\varphi$），因为 $\eta_I$ 定义在 $X$ 上而取值要被搬到 $P$ 的切空间；$F^I_\omega$ 表示 $(F^1_\omega)^{i_1}\cdots(F^N_\omega)^{i_N}$。核对总次数：$\deg[(\dd_\omega\varphi)^*\eta_I]+\deg F^I_\omega = (k-2|I|)+2|I| = k$。]

\begin{theorem}
[18] 的命题 2.2。以下三条性质成立：
\begin{enumerate}
\item $\eta(\varphi,\omega)$ 是不变的且水平的；即
\[ \cL_u\eta(\varphi,\omega) = \iota_u\eta(\varphi,\omega) = 0 \]
对一切垂直向量 $u\in VP=\ker\dd\pi$ 成立。用微分分次代数的语言说，$\eta(\varphi,\omega)$ 是 $P$ 上的基本 $k$-形式（basic $k$-form）。
\item 设 $\rho\in\Omega^{k+1}_G(X)$ 是把 Cartan 微分 $\dd_C$ 作用在 $\eta$ 上所得的等变 $(k+1)$-形式：
\[ \rho = \dd_C\eta = \dd_X\eta - \zeta^a\iota_{V_a}\eta . \]
那么它们各自的拉回 $\eta(\varphi,\omega)$ 与 $\rho(\varphi,\omega)$ 由 $P$ 上普通的外微分相联系：
\[ \rho(\varphi,\omega) = \dd\left(\eta(\varphi,\omega)\right). \]
\item 设映射 $\varphi: P\to X$ “等变同伦”于映射 $\varphi': P\to X$（即 $\varphi$ 可以光滑地形变为 $\varphi'$，使得形变的每一阶段都仍是等变映射）。又设联络 $\omega$ 在联络空间中同伦于 $\omega'$。那么 $\eta(\varphi,\omega)$ 与 $\eta(\varphi',\omega')$ 相差 $P$ 上的一个正合 $k$-形式：
\[ \eta(\varphi',\omega') - \eta(\varphi,\omega) = \dd\Upsilon \]
其中 $\Upsilon\in\Omega^{k-1}(P)$。此外，$\Upsilon$ 也是基本的。
\end{enumerate}
\end{theorem}

% [Tutor Note: p.148 性质 2 中 $\dd_C^2=0$ 的依据是 $\eta\in\Omega^\bullet_G(X)$（即被所有 $\cL_{V_a}$ 零化）：$\dd_C\eta=\dd_X\eta-\zeta^a\iota_{V_a}\eta$，再作用一次 $\dd_C$ 得 $-\zeta^a(\dd_X\iota_{V_a}+\iota_{V_a}\dd_X)\eta=-\zeta^a\cL_{V_a}\eta=0$。这正是前面 $\left[\Omega^\bullet(X)\right]^G$ 这一条件的用处，可用来核对 [18] 中该定理的陈述是否自洽。]

这个定理的含义如下：由性质 1)，形式 $\eta(\varphi,\omega)$ 下降到基流形上的一个 $k$-形式 $\tilde{\eta}(\varphi,\omega)$。（局部下降形式彼此一致，其含义是：在重叠 $U^{(\alpha)}\cap U^{(\beta)}$ 上有 $s^{(\alpha)*}\eta(\varphi,\omega) = s^{(\beta)*}\eta(\varphi,\omega)$；因此它们拼接成 $M$ 上一个整体定义的 $k$-形式。）由性质 2)，如果 $\eta$ 是 $X$ 上 $\dd_C$-闭的等变 $k$-形式，那么 $\tilde{\eta}(\varphi,\omega)$ 是 $M$ 上的闭 $k$-形式（通常意义下）；而如果 $X$ 上两个等变形式相差一个 $\dd_C$-正合项，那么对固定的一对 $(\varphi,\omega)$，它们的拉回相差 $M$ 上的一个 $\dd$-正合项。因此，如果 $S\subset M$

% [Tutor Note: p.148 $s^{(\alpha)}$ 指 $U^{(\alpha)}$ 上的局部截面（原文未另行定义，按上下文即平凡化给出的局部截面 $s^{(\alpha)}: U^{(\alpha)}\to P$）；两个局部下降形式在重叠上相等，靠的正是 $s^{(\alpha)}$ 与 $s^{(\beta)}$ 相差一个转移函数、而 $\eta(\varphi,\omega)$ 是基本形式（性质 1）。]

% [原文 p.149]
是底流形的一个 $k$ 维、定向、无边的子流形，那么由积分
\[ S(\eta,\varphi,\omega) \equiv \int_S \tilde{\eta}(\varphi,\omega) \]
所得的数目与 $\eta$ 的等变上同调类的代表无关。

% [Tutor Note: p.149 原文末句 “In case $k=m$ and $S=M$” 有笔误：$M$ 是 $2m$ 维的（见 p.142），故 $S=M$ 时须 $k=\dim M=2m$ 而非 $k=m$；否则 $\tilde{\eta}$ 的次数与积分维数不符。正确形式应为 $k=2m$ 且 $S=M$，此时 $\int_M\tilde{\eta}(\varphi,\omega)$ 即规范 $\sigma$-模型作用量中的拓扑项。]

由性质 3)，如果允许 $\varphi$ 与 $\omega$ 在截面与联络空间的一个连通分支上连续变化，那么数 $S(\eta,\varphi,\omega)$ 不变；我们把它写作
\[ S(\eta,\varphi,\omega) = S([\eta],[\varphi],[\omega]) \]
以表明它只依赖于 $\eta$ 的等变上同调类以及 $\varphi$ 与 $\omega$ 的同伦类；正因如此，我们说这个数是“拓扑的”。当 $k=m$ 且 $S=M$ 时，这就给出一个可以纳入规范 $\sigma$-模型作用量的拓扑项。
